\chapter{GPT and text generation}
\label{ch:gpt-generation}

\begin{quote}
\emph{``GPT is just next-token prediction, scaled up. The emergent capabilities come from the scale, not from a change in the objective.''}
\end{quote}

% ============================================================
\section{The GPT family}

GPT (Generative Pre-trained Transformer) models are \emph{decoder-only} Transformers trained with a causal language modeling objective: predict the next token given all previous tokens.

\begin{longtable}{p{2.5cm}p{2.5cm}p{2.5cm}p{4.5cm}}
\toprule
\textbf{Model} & \textbf{Parameters} & \textbf{Year} & \textbf{Key innovation} \\
\midrule
GPT-1 & 117M & 2018 & Pre-train + fine-tune paradigm \\
GPT-2 & 1.5B & 2019 & Zero-shot task transfer \\
GPT-3 & 175B & 2020 & In-context learning, few-shot \\
GPT-4 & Unknown & 2023 & Multimodal, RLHF \\
GPT-4o & Unknown & 2024 & Omni: text, vision, audio \\
\bottomrule
\end{longtable}

For hands-on work, we use \textbf{GPT-2} (open weights, runs on free Colab GPU) and \textbf{TinyLlama-1.1B} (modern architecture, small enough for experimentation).

% ============================================================
\section{Autoregressive generation}

Autoregressive generation works token by token:
\begin{enumerate}
  \item Feed the prompt tokens to the model.
  \item The model outputs logits over the entire vocabulary.
  \item Select the next token from the logits (via a \emph{decoding strategy}).
  \item Append the token to the sequence. Repeat.
\end{enumerate}

\begin{minted}{python}
from transformers import GPT2LMHeadModel, GPT2Tokenizer
import torch

tokenizer = GPT2Tokenizer.from_pretrained("gpt2")
model = GPT2LMHeadModel.from_pretrained("gpt2")
model.eval()

prompt = "The future of artificial intelligence"
input_ids = tokenizer.encode(prompt, return_tensors="pt")

# Manual autoregressive loop
generated = input_ids.clone()
for _ in range(30):
    with torch.no_grad():
        outputs = model(generated)
        next_logits = outputs.logits[:, -1, :]  # logits for last position
        next_token = torch.argmax(next_logits, dim=-1, keepdim=True)  # greedy
        generated = torch.cat([generated, next_token], dim=-1)

print(tokenizer.decode(generated[0]))
\end{minted}

% ============================================================
\section{Decoding strategies}

\subsection{Greedy decoding}

Always pick the highest-probability token. Fast but repetitive and boring.

\subsection{Temperature sampling}

Scale the logits by a temperature $T$ before softmax:
\[
P(w_i) = \frac{e^{z_i / T}}{\sum_j e^{z_j / T}}
\]
\begin{itemize}
  \item $T = 1.0$: standard sampling.
  \item $T < 1.0$: sharper distribution, more deterministic.
  \item $T > 1.0$: flatter distribution, more creative/random.
\end{itemize}

\subsection{Top-k sampling}

Keep only the $k$ highest-probability tokens, redistribute probability among them.

\subsection{Top-p (nucleus) sampling}

Keep the smallest set of tokens whose cumulative probability exceeds $p$. Adapts dynamically: uses fewer tokens when the model is confident, more when uncertain.

\subsection{Beam search}

Maintain $b$ candidate sequences (beams) at each step. Select the $b$ highest-probability continuations. Good for translation and summarization where quality matters more than diversity.

\begin{minted}{python}
from transformers import pipeline

generator = pipeline("text-generation", model="gpt2")

prompt = "In 2026, the most important AI breakthrough was"

# Greedy
print("=== Greedy ===")
print(generator(prompt, max_new_tokens=50, do_sample=False)[0]["generated_text"])

# Temperature sampling
print("\n=== Temperature 0.3 (focused) ===")
print(generator(prompt, max_new_tokens=50, do_sample=True,
                temperature=0.3)[0]["generated_text"])

print("\n=== Temperature 1.5 (creative) ===")
print(generator(prompt, max_new_tokens=50, do_sample=True,
                temperature=1.5)[0]["generated_text"])

# Top-k sampling
print("\n=== Top-k (k=10) ===")
print(generator(prompt, max_new_tokens=50, do_sample=True,
                top_k=10)[0]["generated_text"])

# Top-p (nucleus) sampling
print("\n=== Top-p (p=0.9) ===")
print(generator(prompt, max_new_tokens=50, do_sample=True,
                top_p=0.9)[0]["generated_text"])

# Beam search
print("\n=== Beam search (4 beams) ===")
print(generator(prompt, max_new_tokens=50, num_beams=4,
                do_sample=False)[0]["generated_text"])
\end{minted}

\begin{aitip}
For most creative tasks, \texttt{temperature=0.7} with \texttt{top\_p=0.9} is a good starting point. For factual tasks (code, math), use \texttt{temperature=0.2} or greedy decoding.
\end{aitip}

% ============================================================
\section{HuggingFace \texttt{generate()} API}

The \texttt{generate()} method on any HuggingFace causal LM supports all decoding strategies:

\begin{minted}{python}
from transformers import AutoModelForCausalLM, AutoTokenizer

model_name = "TinyLlama/TinyLlama-1.1B-Chat-v1.0"
tokenizer = AutoTokenizer.from_pretrained(model_name)
model = AutoModelForCausalLM.from_pretrained(
    model_name, torch_dtype="auto", device_map="auto"
)

messages = [
    {"role": "system", "content": "You are a helpful AI assistant."},
    {"role": "user", "content": "Explain gradient descent in 3 sentences."},
]
prompt = tokenizer.apply_chat_template(messages, tokenize=False,
                                        add_generation_prompt=True)
inputs = tokenizer(prompt, return_tensors="pt").to(model.device)

output = model.generate(
    **inputs,
    max_new_tokens=150,
    temperature=0.7,
    top_p=0.9,
    do_sample=True,
    repetition_penalty=1.1,
)
print(tokenizer.decode(output[0], skip_special_tokens=True))
\end{minted}

% ============================================================
\section{Controlling repetition}

Base models tend to repeat themselves. Strategies:

\begin{minted}{python}
# Repetition penalty: penalizes tokens that already appeared
output = model.generate(**inputs, max_new_tokens=100,
                        repetition_penalty=1.2)

# No-repeat n-gram: prevents repeating any n-gram
output = model.generate(**inputs, max_new_tokens=100,
                        no_repeat_ngram_size=3)
\end{minted}

\begin{warningbox}
Setting \texttt{repetition\_penalty} too high (> 1.5) can cause incoherent output because the model avoids common, necessary words. Start with 1.1--1.2.
\end{warningbox}

% ============================================================
\section{Comparing models: GPT-2 vs TinyLlama}

\begin{minted}{python}
from transformers import pipeline

models = ["gpt2", "TinyLlama/TinyLlama-1.1B-Chat-v1.0"]
prompt = "Machine learning is"

for m in models:
    gen = pipeline("text-generation", model=m, device_map="auto")
    result = gen(prompt, max_new_tokens=60, do_sample=True,
                 temperature=0.7, top_p=0.9)
    print(f"\n=== {m} ===")
    print(result[0]["generated_text"])
\end{minted}

\begin{exercise}
\begin{enumerate}
  \item Generate 5 completions for the same prompt using temperature values 0.1, 0.5, 0.8, 1.0, and 1.5. Describe how the outputs change.
  \item Implement top-k sampling from scratch: given logits, zero out everything except the top $k$ values, apply softmax, and sample.
  \item Use TinyLlama to generate a 200-word explanation of photosynthesis. Try greedy, beam search ($b=5$), and nucleus sampling ($p=0.9$). Which is best?
  \item Measure the generation speed (tokens/second) of GPT-2 on CPU vs GPU.
  \item Generate a short story with GPT-2. First without repetition penalty, then with \texttt{repetition\_penalty=1.15}. Compare.
\end{enumerate}
\end{exercise}

\begin{ethicsbox}
Open-weight models like GPT-2 and TinyLlama have no built-in content filters. They can generate hate speech, disinformation, and harmful instructions. When building applications, always add output filtering, content moderation, or use instruction-tuned models with safety training.
\end{ethicsbox}

% ============================================================
\section{Chapter summary}

\begin{itemize}
  \item GPT models are decoder-only Transformers trained to predict the next token.
  \item Autoregressive generation produces text one token at a time, feeding each output back as input.
  \item Decoding strategies (greedy, temperature, top-k, top-p, beam search) control the tradeoff between quality and diversity.
  \item Temperature scales logits: low = focused, high = creative.
  \item Top-p (nucleus) sampling adapts the candidate set dynamically to the model's confidence.
  \item HuggingFace \texttt{generate()} supports all strategies with simple keyword arguments.
  \item Repetition penalties prevent degenerate loops in open-ended generation.
\end{itemize}
